\begin{enumerate}
\item \textbf{O que foi refeito.} Só a medição, sob o protocolo comum v1: imagem 4.1.12 por digest, container dedicado, os 9 deployments admissíveis + baseline + 2 réplicas A/A, 5 repetições a frio, 60 execuções aceitas (nenhuma descartada) e \envrawreads{} leituras conferidas pelo oráculo independente. Todas as execuções passaram pela auditoria integral de registros.
\item \textbf{Leitura.} Há dois grupos bem separados. Quem lê de uma projeção TS3 (partição por sensor) fica em 2,91--2,96\,ms, com 2--3 requisições CQL. Quem lê direto de TS0/TS1 fica em 4,57--5,06\,ms, com 4--6 requisições. A razão é $\approx$0,63, muito acima do piso de ruído (6,1\%). Dentro de cada grupo as diferenças são empate.
\item \textbf{Ingestão e disco.} Deployments de 1 tabela ingerem $\approx$380\,ev/s; os de 2 tabelas, $\approx$243\,ev/s (0,64$\times$). Os de 2 tabelas também ocupam $\approx$2,1$\times$ o disco. Há, portanto, um compromisso real entre leitura e escrita/armazenamento.
\item \textbf{Preservação.} O melhor medido (TS1m+TS3m, HBench \#3) e todos os empatados com ele estão no HBench: Top-3 sim, near-best sim, regret 0. O HBench retém 100\% do Pareto 2D (leitura $\times$ ingestão, critério da Etapa 9.1) e 83\% do Pareto 3D, porque fica de fora TS1h+TS3m, que empata com os demais. Em compensação, o escore do seletor não prediz a latência: $\tau=-0{,}06$, e o deployment de menor escore (TS0, o controle) é dos mais lentos. Só com $k=1$ o regret seria de 61\%.
\item \textbf{Baseline TPC.} \texttt{b\_tpcx\_logical\_raw} (uma tabela particionada por \texttt{sensor\_name}) empata em leitura com o melhor deployment do SchemaLens (razões 1,02--1,04, empate). Ele ingere 1,6$\times$ mais rápido que os deployments TS3, que precisam de 2 tabelas, e ocupa metade do disco. Nesta carga e nesta escala o baseline domina os deployments TS3 em escrita e armazenamento, sem perder em leitura.
\item \textbf{Ablações.} A regra decisiva é a de TS3: sem ela o regret sobe a 57\%. Sem TS0 o Stage 8 fica inviável, porque o controle é obrigatório. Sem TS1 o regret é de 0,2\% (empate). M1 e M3 não mudam os candidatos.
\end{enumerate}
